\documentclass[11pt]{article}
\usepackage[utf8]{inputenc}
\usepackage[spanish]{babel}
\usepackage{geometry}
\geometry{a4paper, margin=1in}

% Fuentes y Estética
\usepackage{helvet}
\renewcommand{\familydefault}{\sfdefault}
\usepackage[table]{xcolor}
\usepackage{tcolorbox}

% Colores Corporativos
\definecolor{IndustrialBlue}{RGB}{0, 51, 102}

\title{Especificaciones Técnicas: Instrumentación Lutron}
\author{Localización Técnica Profesional}
\date{Febrero 2026}

\begin{document}
\maketitle

\section{Introducción}
Este documento es una versión optimizada del manual de fábrica del modelo **8007fe**, rediseñada para mejorar la legibilidad y la precisión técnica mediante tipografía profesional.

\section{Especificaciones Eléctricas}
\begin{center}
{\renewcommand{\arraystretch}{1.5} % Más espacio entre filas para legibilidad
\begin{tabular}{|>{\columncolor{IndustrialBlue!10}\bfseries}l|l|l|}
\hline
\rowcolor{IndustrialBlue}
\textcolor{white}{Rango} & \textcolor{white}{Resolución} & \textcolor{white}{Precisión} \\ \hline
200.0 mT & 0.1 mT & $\pm$(5\% + 3d) \\ \hline
2,000 mT & 1 mT & $\pm$(5\% + 3d) \\ \hline
\end{tabular}}
\end{center}

\section{Notas de Seguridad}
\begin{itemize}
    \item No exceder el límite máximo de entrada del sensor.
    \item Utilizar baterías de 9V certificadas para garantizar la estabilidad de la lectura.
\end{itemize}

\end{document}